\documentclass[10pt,a4paper]{amsart}
\usepackage[margin=19mm]{geometry}
\usepackage{amsmath,amssymb,mathtools}
\usepackage[scr=rsfs,cal=euler]{mathalfa}
\usepackage{xfrac}
\usepackage[unicode,hidelinks,bookmarks=false]{hyperref}
\newcommand{\ie}{i.e.\ }
\newcommand{\cf}{cf.\ }
\newcommand{\resp}{respectively\ }
\newcommand{\Subsection}{\subsection}
\providecommand{\nobreakdash}{\nobreak\hbox{-}}
\setlength{\emergencystretch}{2em}
\multlinegap0pt
\usepackage{bm}
\usepackage{bbm}
\usepackage{dsfont}
\usepackage{xspace}
\usepackage{cancel}
\usepackage{mathtools}
\usepackage{amsthm}
\makeatletter
\@ifundefined{theorem}{\newtheorem{theorem}{Theorem}}{}
\@ifundefined{lemma}{\newtheorem{lemma}[theorem]{Lemma}}{}
\@ifundefined{proposition}{\newtheorem{proposition}[theorem]{Proposition}}{}
\makeatother
\def\de{\delta}
\let\f=\frac
\title[Enhanced dissipation and Taylor dispersion by a parallel she]{Enhanced dissipation and Taylor dispersion by a parallel shear flow in an infinite cylinder with unbounded cross section}
\author{Te Li, Le Zhang}
\date{Source: arXiv:2510.13097v1 (CC BY 4.0)}
\begin{document}
\maketitle

\noindent\emph{Source and licence.} Enhanced dissipation and Taylor dispersion by a parallel shear flow in an infinite cylinder with unbounded cross section. Version arXiv:2510.13097v1 (CC BY 4.0), distributed under the Creative Commons Attribution 4.0 International licence, as recorded in the arXiv licence metadata for this exact version and shown on the abstract page https://arxiv.org/abs/2510.13097v1. Full rights notice: licenses/arxiv-2510.13097-CC-BY-4.0.txt.\par\smallskip

\noindent\textit{Editorial scope.} Counted unit: the Proposition whose source label is \texttt{mathcal e measure-prop} in the article (source0.tex lines 1964--1970, source label \texttt{mathcal e measure-prop}), delivered together with the article's own complete original proof of it (source0.tex lines 1971--2157, one \texttt{proof} environment of 187 lines). No mathematics is written, repaired or re-derived: statement and proof are the article's own text line for line. Required context retained uncounted: eq:geq (lines 273--277); thm:main1 (lines 299--327); eq:vder (lines 302--304); non-degenerate at infinity (lines 306--308); mathcal{E (lines 1792--1802). Every definition, display and label of those lines is retained unchanged. Omitted from this package and not redistributed: 1--272, 278--298, 305--305, 309--1791, 1803--1963, 2158--2372. Nothing inside a retained excerpt is omitted or altered: every retained body is delivered exactly as the article's lines are written, so no omission receipt of any kind accompanies this package. Citations: the retained excerpts carry no citation command at all, so no bibliography entry of the article is transcribed and no reference is redistributed. Reference adaptation: none was needed and none was made; every reference inside the delivered unit resolves inside it, so no unresolved reference or citation is produced. Printed slips kept literal: no printed word, symbol, hypothesis or number of the counted statement or of its proof was altered. Layout and class: the article's own class is \texttt{amsart} with packages including amssymb, amsmath, latexsym, amsfonts, amsbsy, bbm, amsthm, mathtools, graphicx, CJKutf8, CJKnumb, CJKulem, color, xcolor; the wrapper is the lane's pinned \texttt{amsart} editorial wrapper at 10pt on A4, which restates the theorem-like environments the retained text needs and adds the article's own macro definitions that the retained text needs (\texttt{\textbackslash de}, \texttt{\textbackslash f}), with extra wrapper packages (bm, bbm, dsfont, xspace, cancel, mathtools). The delivered numbering is the unit's own. Assets: no retained span contains a figure environment, a graphics-inclusion command, a PDF-page inclusion, a table or a TikZ picture; no image file is copied into this package, the package declares no asset dependency and no image file is redistributed with it. Licence: version arXiv:2510.13097v1 of this article is distributed under the Creative Commons Attribution 4.0 International licence, as recorded in the arXiv licence metadata for this exact version and shown on the abstract page https://arxiv.org/abs/2510.13097v1. A full rights notice is delivered at licenses/arxiv-2510.13097-CC-BY-4.0.txt. Characters: every retained excerpt is pure ASCII, so the delivered engine, XeLaTeX with its default Latin Modern text font, reports no missing character.
\begin{equation}\label{eq:geq}
  \partial_t g(k,y,t) + \mathrm{i}k\,v(y)\,g(k,y,t)
  \;=\; \nu\,\Delta_y g(k,y,t)\,,
  \qquad y\in\Omega,\; t>0,
\end{equation}

\begin{theorem}[One-dimensional case]\label{thm:main1}
Assume that $d = 1, m\in\mathbb{N}^{*}, v : \overline{\Omega} \to \mathbb{R}$ belongs to $C^m(\overline{\Omega})$ and satisfies the following conditions:  
\begin{enumerate}
    \item[(1)] The derivatives of $v$ up to order $m$ do not vanish simultaneously, namely: \begin{equation}\label{eq:vder} |v'(y)| + |v''(y)| + \cdots + |v^{(m)}(y)| > 0,
\qquad \text{for all } y \in \overline{\Omega}.
\end{equation} 
    \item[(2)] $v$ is non-degenerate at infinity, namely: 
  \begin{equation}\label{non-degenerate at infinity}
\exists\, c_0>0 \quad\text{s.t.} \quad  \varliminf_{|y|\to\infty}|v'(y)|\geq c_0>0.
\end{equation} 
\end{enumerate}
Here, $\Omega$ is taken to be either $\mathbb{R}$, $(-\infty, L_1)$, or $(L_2, \infty)$, where $L_1, L_2$ are arbitrary real numbers. 

Then, there exist positive constants $C$ and $c$ such that, for every $\nu > 0$, 
every $k \neq 0$, and every initial datum $g_0 \in L^2(\Omega)$, the solution 
of \eqref{eq:geq} satisfies, for all $t \ge 0$, 
\begin{equation}\label{lamnuk}
  \|g(k,t)\|_{L^2(\Omega)} \,\le\, 
  C\,\mathrm{e}^{-c \lambda_{\nu,k} t}\,\|g_0\|_{L^2(\Omega)}\,, 
\end{equation}
where 
\begin{equation*}
  \lambda_{\nu,k} \,=\,
  \begin{cases} 
 \nu^{\tfrac{m}{m+2}} |k|^{\tfrac{2}{m+2}}, & \text{if }~ 0 < \nu \le |k|\,, \\[1ex]
  \dfrac{k^2}{\nu}, & \text{if }~ 0 < |k| \le \nu\,.
  \end{cases}
\end{equation*}
\end{theorem}

\begin{lemma}\label{mathcal{E}(lambda)}
Let $d=1$, $m \geq 1$, and $L_1, L_2$ are arbitrary real numbers with $L_1 < L_2$. Suppose $v:[L_1,L_2] \to \mathbb{R}$ is a $C^m$ function satisfying the $m$-th order non-degeneracy condition
\[
 |v'(y)| + |v''(y)| + \cdots + |v^{(m)}(y)| > 0,\quad \forall\, y \in [L_1,L_2].
\]
Then there exist constants $C_0' > 0$ and $\hat{\delta}_0 \in (0,1)$ such that, 
for any $\delta \in (0,\hat{\delta}_0]$ and any $\lambda \in \mathbb{R}$, we have
\[
 m(\mathcal{E}^m_{\lambda,\delta})\leq C_0' \delta .
\]
\end{lemma}

\begin{proposition}\label{mathcal e measure-prop}
Assume $d=1$ and $v : \Omega\to \mathbb{R}$ is a $C^{m}$ function, where $m$ and $\Omega$ are the same as in Theorem \ref{thm:main1}, and $v$ satisfies conditions \eqref{eq:vder} and \eqref{non-degenerate at infinity} of Theorem \ref{thm:main1}. Then
there exist constant $C > 0$ and  $\delta_0>0$ such that, for all $\nu > 0$, all $k \neq 0$ and all $\lambda\in\mathbb{R}\ and\ \delta\in(0,\delta_0],$ 
\begin{equation}\label{mathcal e measure}
m(\mathcal{E}_{\lambda,\delta}^m) \leq C\delta.
\end{equation}
\end{proposition}

\begin{proof}[\bf Proof of Proposition \ref{mathcal e measure-prop}]
Without loss of generality, we present the proof for $\Omega=\mathbb{R}$.

Since $\varliminf_{|y|\to\infty}|v'(y)|\ge c_0>0$, there exists $\Delta>0$ with
\begin{equation}\label{xjx}
  |v'(y)| \ge \tfrac{c_0}{2} \quad \text{for all } |y|\ge \Delta.
\end{equation}

We now prove that $m(\mathcal{E}^{m}_{\lambda,\delta}) = O(\delta)$ uniformly for $\lambda\in\mathbb{R}$; that is, there exist constants $C_0>0$ and $\delta_0>0$ such that, for any $\lambda\in\mathbb{R}$ and any $\delta\in(0,\delta_0]$,
\[
  m(\mathcal{E}^{m}_{\lambda,\delta}) \le C_0\,\delta.
\]

From \eqref{xjx}, either $ v'(y)\ge \tfrac{c_0}{2} \text{ or } v'(y)\le -\tfrac{c_0}{2} (\forall y\in[\Delta,\infty))$. Without loss of generality, assume that $v'(y)\ge \tfrac{c_0}{2}$ holds. Then $v$ is strictly increasing on $[\Delta,\infty)$, and for all $y\ge \Delta$,
\[
  v(y)=v(\Delta)+\int_{\Delta}^{y} v'(z)\,dz
  \;\ge\; v(\Delta)+\frac{c_0}{2}\,(y-\Delta),
\]
hence $v(+\infty)=+\infty$.

For any $\delta\in(0,\Delta]$ and any $\lambda\in\mathbb{R}$, set
\begin{align}\label{FF}
\begin{split}
  E^{m}_{\lambda,\delta}
  = &\bigl(E^{m}_{\lambda,\delta}\cap(-\infty,-2\Delta]\bigr)
     \;\cup\; \bigl(E^{m}_{\lambda,\delta}\cap[-2\Delta,2\Delta]\bigr)
     \;\cup\; \bigl(E^{m}_{\lambda,\delta}\cap[2\Delta,\infty)\bigr) \\
   =&:F_1 \cup F_2 \cup F_3 .
\end{split}
\end{align}
Define
\begin{align}\label{G}
\begin{split}
 & \mathcal{E}^{m}_{\lambda,\delta}
  = \{y\in\mathbb{R}:\mathrm{dist}(y,E^{m}_{\lambda,\delta})<\delta\} \\
  = &\{y\in\mathbb{R}:\mathrm{dist}(y,F_1)<\delta\}
     \;\cup\; \{y\in\mathbb{R}:\mathrm{dist}(y,F_2)<\delta\}
     \;\cup\; \{y\in\mathbb{R}:\mathrm{dist}(y,F_3)<\delta\} \\
 =&: G_1 \cup G_2 \cup G_3 .
\end{split}
\end{align}
Moreover, define
\begin{align}\label{hh}
\begin{split}
  H_1 &:= E^{m}_{\lambda,\delta}\cap(-\infty,-\Delta)
         \;\supset\; E^{m}_{\lambda,\delta}\cap(-\infty,-2\Delta] \;=\; F_1, \\
  H_2 &:= E^{m}_{\lambda,\delta}\cap(-3\Delta,3\Delta)
         \;\supset\; E^{m}_{\lambda,\delta}\cap[-2\Delta,2\Delta] \;=\; F_2, \\
  H_3 &:= E^{m}_{\lambda,\delta}\cap(\Delta,\infty)
         \;\supset\; E^{m}_{\lambda,\delta}\cap[2\Delta,\infty) \;=\; F_3 .
\end{split}
\end{align}
Then
\begin{align}\label{GG}
\begin{split}
  &G_1 \subset (-\infty,-2\Delta+\delta) \subset (-\infty,-\Delta),\\
  &G_2 \subset (-2\Delta-\delta,\,2\Delta+\delta) \subset (-3\Delta,3\Delta),\\
  &G_3 \subset (2\Delta-\delta,\infty) \subset (\Delta,\infty).
\end{split}
\end{align}
Consequently,
\begin{align}\label{gs}
\begin{split}
  G_3
  &= \{y\in\mathbb{R}:\mathrm{dist}(y,F_3)<\delta\} \overset{\eqref{GG}}{=} \{y\in(\Delta,\infty):\mathrm{dist}(y,F_3)<\delta\} \\
  &\overset{\eqref{hh}}{\subset} \{y\in(\Delta,\infty):\mathrm{dist}(y,H_3)<\delta\}.
\end{split}
\end{align}

We distinguish the following three cases:
\begin{align}\label{sz}
\begin{split}
  &\text{Case e.1: } v(\Delta)\le \lambda-\delta,\\
  &\text{Case e.2: } \lambda-\delta< v(\Delta)< \lambda+\delta,\\
  &\text{Case e.3: } \lambda+\delta\le v(\Delta).
\end{split}
\end{align}

\begin{itemize}
    \item For Case e.1, by the strict monotonicity of $v$,
there exists a unique pair $(\Delta_1,\Delta_2)$ with $\Delta_1\ge \Delta$ and
$\Delta_2\ge \Delta$ such that
\begin{equation}\label{dd}
  v(\Delta_1)=\lambda-\delta,\qquad v(\Delta_2)=\lambda+\delta,
\end{equation}
and moreover
\begin{equation}\label{jz}
  \Delta_1<\Delta_2,\qquad H_3=(\Delta_1,\Delta_2).
\end{equation}
Hence, by \eqref{gs} and \eqref{jz},
\[
  G_3 \subset \{\,y\in(\Delta,\infty): \text{dist}\bigl(y,(\Delta_1,\Delta_2)\bigr)<\delta \,\}
      \subset (\Delta_1-\delta,\Delta_2+\delta),
\]
and thus
\begin{equation}\label{dede}
  m(G_3) \le (\Delta_2+\delta)-(\Delta_1-\delta) = 2\delta + \Delta_2-\Delta_1.
\end{equation}
Moreover,
\begin{align*}
  2\delta
  &= (\lambda+\delta)-(\lambda-\delta)
   \overset{\eqref{dd}}{=} v(\Delta_2)-v(\Delta_1)
   = \int_{\Delta_1}^{\Delta_2} v'(y)\,dy
   \ge \int_{\Delta_1}^{\Delta_2} \frac{c_0}{2}\,dy
   = \frac{c_0}{2}\,(\Delta_2-\Delta_1),
\end{align*}
whence
\begin{equation}\label{ddd}
  \Delta_2-\Delta_1 \le \frac{4}{c_0}\,\delta.
\end{equation}
Combining \eqref{dede} and \eqref{ddd}, we obtain
\begin{equation}\label{ggg}
  m(G_3) \le \Bigl(2+\frac{4}{c_0}\Bigr)\delta.
\end{equation}
\item For Case e.2, by the strict monotonicity of $v$,
there exists a unique $\Delta_3>\Delta$ such that
\[
  v(\Delta_3)=\lambda+\delta,
  \qquad H_3=(\Delta,\Delta_3).
\]
Hence, by \eqref{gs} and \eqref{jz},
\[
  G_3 \subset \{\,y\in(\Delta,\infty): \text{dist}\bigl(y,(\Delta,\Delta_3)\bigr)<\delta \,\}
      \subset (\Delta-\delta,\Delta_3+\delta),
\]
and thus
\begin{equation}\label{dedede}
  m(G_3) \le (\Delta_3+\delta)-(\Delta-\delta) = 2\delta + \Delta_3-\Delta.
\end{equation}
Moreover, using $v(\Delta_3)=\lambda+\delta$ and $\lambda-\delta< v(\Delta)$,
\[
  v(\Delta_3)-v(\Delta) \le (\lambda+\delta)-(\lambda-\delta)=2\delta,
\]
then 
\[
  v(\Delta_3)-v(\Delta)=\int_{\Delta}^{\Delta_3} v'(y)\,dy
  \ge \int_{\Delta}^{\Delta_3} \frac{c_0}{2}\,dy
  = \frac{c_0}{2}\,(\Delta_3-\Delta).
\]
Hence
\begin{equation}\label{dxx}
  \Delta_3-\Delta \le \frac{4}{c_0}\,\delta.
\end{equation}
Combining \eqref{dedede} and \eqref{dxx}, we obtain
\begin{equation}\label{gsgs}
  m(G_3) \le \Bigl(2+\frac{4}{c_0}\Bigr)\delta.
\end{equation}
\item For Case e.3, by the monotonicity of $v$, since for all $y\ge 2\Delta$ we have
\( v(y) > v(\Delta) \ge \lambda+\delta \),
it follows that \(F_3=\varnothing\). Hence, by \eqref{G},
\begin{equation}\label{lin}
  G_3=\varnothing,\qquad m(G_3)=0.
\end{equation}

\end{itemize}

Therefore, combining \eqref{ggg}, \eqref{gsgs}, and \eqref{lin}, for any $ \delta\in(0,\Delta]$ and $\lambda\in\mathbb{R}$, regardless of which case in \eqref{sz} occurs, we have
\begin{equation}\label{gggg}
  m(G_3)\le \Bigl(2+\frac{4}{c_0}\Bigr)\delta.
\end{equation}

Similarly, for any $\delta\in(0,\Delta]$ and any $\lambda\in\mathbb{R}$,
\begin{equation}\label{ggggg}
  m(G_1)\le \Bigl(2+\frac{4}{c_0}\Bigr)\delta.
\end{equation}

For any $y\in[-3\Delta,3\Delta]\subset\mathbb{R}$, the $m$-th order non-degeneracy condition holds:
\[
  |v'(y)|+\cdots+|v^{(m)}(y)|>0.
\]
By Lemma \ref{mathcal{E}(lambda)}, there exist constants $C_0'>0$ and $\hat{\delta}_0>0$ such that,
for any $\delta\in(0,\hat{\delta}_0]$ and any $\lambda\in\mathbb{R}$,
\begin{equation}\label{gg}
  m(G_2)\le C_0'\,\delta .
\end{equation}

Therefore, combining \eqref{lin}, \eqref{ggggg}, and \eqref{gg}, by taking $$\delta_0=\min\{\Delta,\hat{\delta}_0\},\quad C=4+\f{8}{c_0}+C_0',$$ we obtain that for any $\delta\in(0,\delta_0]$ and any $\lambda\in\mathbb{R},$
\begin{align*}
  m(\mathcal{E}^{m}_{\lambda,\delta})
 \le m(G_1)+m(G_2)+m(G_3)
  \le \Bigl(4+\frac{8}{c_0}+C_0'\Bigr)\delta=C\de.
\end{align*}


This completes the proof of Proposition \ref{mathcal e measure-prop}.
\end{proof}

\end{document}
